% =====================================================================
% GA-SUM-02 | EXP-1: Técnicas de elicitación y herramientas de apoyo a la IR
% Asignatura: Ingeniería de Requerimientos (ISR-401) - UTEQ
% Compilar con:  pdflatex main  ->  bibtex main  ->  pdflatex main (x2)
% Diseño: minimalista, morado y negro (sin tema Beamer predefinido)
% =====================================================================
\documentclass[aspectratio=169,11pt,t]{beamer}

\usepackage[utf8]{inputenc}
\usepackage[T1]{fontenc}
\usepackage[spanish,es-noshorthands]{babel}
\usepackage{lmodern}
\usepackage{booktabs}
\usepackage{array}
\usepackage{colortbl}
\usepackage{tikz}
\usetikzlibrary{positioning,arrows.meta}
\usepackage{listings}
\usepackage{graphicx}
\usepackage{xcolor}
\usepackage{hyperref}

% ---------------------- DATOS DEL ESTUDIANTE (EDITAR) ----------------
\newcommand{\Estudiante}{[Rivera Rodriguez Joseph Jeremy]}
\newcommand{\Correo}{[jriverar4@uteq.edu.ec]}
\newcommand{\Semestre}{[Cuarto Semestre / Paralelo A]}
\newcommand{\Carrera}{[Software]}
\newcommand{\Equipo}{FNR}
\newcommand{\Docente}{Dr. Gleiston Cicerón Guerrero Ulloa}
\newcommand{\Fecha}{21 de noviembre de 2026}
\newcommand{\RepoURL}{https://github.com/jriverar4-svg/ISR401--FNR-.git}

% ---------------------- PALETA: MORADO Y NEGRO -----------------------
\definecolor{tinta}{HTML}{15121C}        % negro con matiz violeta
\definecolor{acento}{HTML}{5B34B8}       % morado (sobre fondo claro)
\definecolor{acentoclaro}{HTML}{B79CFF}  % morado claro (sobre fondo oscuro)
\definecolor{lila}{HTML}{F1EDFA}         % tinte muy suave
\definecolor{gris}{HTML}{6E6A7A}
\definecolor{codigofondo}{HTML}{15121C}
\definecolor{codigotexto}{HTML}{E9E5F5}

% ---------------------- ESTILO GLOBAL --------------------------------
\setbeamercolor{normal text}{fg=tinta,bg=white}
\setbeamercolor{alerted text}{fg=acento}
\setbeamercolor{structure}{fg=acento}
\setbeamersize{text margin left=1.4cm,text margin right=1.4cm}
\setbeamertemplate{navigation symbols}{}
\setbeamertemplate{itemize item}{\color{acento}\tiny$\blacksquare$}
\setbeamertemplate{itemize subitem}{\color{acento}\tiny$\blacksquare$}
\setbeamertemplate{enumerate item}{\color{acento}\bfseries\insertenumlabel.}
\setbeamertemplate{bibliography item}[text]
\setbeamercolor{bibliography item}{fg=acento}
\setlength{\leftmargini}{1.1em}
\setbeamertemplate{itemize/enumerate body begin}{\setlength{\itemsep}{3pt}}
\hypersetup{colorlinks=true,urlcolor=acento,linkcolor=acento,citecolor=acento}

% Título de diapositiva: etiqueta morada pequeña + título negro + trazo corto
\setbeamertemplate{frametitle}{%
  \vskip -0.2em
  {\color{acento}\scriptsize\bfseries\MakeUppercase{\insertframesubtitle}}\par
  {\color{tinta}\large\bfseries\insertframetitle\par}
  \vskip 0.2em
  {\color{acento}\rule{1.0cm}{1.4pt}}\par
  \vskip 0.2em}

% Pie de página discreto
\setbeamertemplate{footline}{%
  \leavevmode\hbox to\paperwidth{\hspace{1.4cm}%
    {\color{gris}\tiny Pauly Creative's Craft \ · \ GA-SUM-02 \ · \ EXP-1}\hfill
    {\color{gris}\tiny\insertframenumber\ / \inserttotalframenumber}\hspace{1.4cm}}%
  \vskip 0.7em}

% ---------------------- COMANDOS DE ESTRUCTURA -----------------------
\newcommand{\etiqueta}[1]{{\color{acento}\scriptsize\bfseries\MakeUppercase{#1}}\par\vspace{3pt}}
% Nota con filete lateral morado (sustituye a los cuadros de colores)
\newcommand{\nota}[2]{\par\vspace{3pt}\noindent
  \begin{tabular}{@{}!{\color{acento}\vrule width 2.2pt}@{\hspace{8pt}}p{\dimexpr\linewidth-14pt\relax}@{}}
    {\color{acento}\scriptsize\bfseries\MakeUppercase{#1}}\par\vspace{1pt}#2
  \end{tabular}\par}
% Paso numerado en fila horizontal
\newcommand{\paso}[2]{\parbox[t]{0.178\textwidth}{\raggedright\color{acento}\large\bfseries #1\par\vspace{2pt}\color{tinta}\scriptsize #2}}
% Conclusión numerada
\newcommand{\conc}[2]{{\color{acento}\large\bfseries #1}\par\vspace{2pt}{\small #2}\par\vspace{9pt}}

% ---------------------- LISTINGS (.sdoc) sobre fondo oscuro ----------
\lstdefinestyle{sdoc}{
  basicstyle=\ttfamily\fontsize{7}{8.3}\selectfont\color{codigotexto},
  backgroundcolor=\color{codigofondo},
  frame=none, xleftmargin=6pt, xrightmargin=6pt,
  framexleftmargin=6pt, framexrightmargin=6pt,
  framextopmargin=5pt, framexbottommargin=5pt,
  columns=flexible, keepspaces=true, showstringspaces=false, breaklines=true,
  keywordstyle=\bfseries\color{acentoclaro},
  morekeywords={UID,TITLE,TAGS,STATEMENT,RELATIONS,TYPE,VALUE},
  moredelim=[s][\bfseries\color{acentoclaro}]{[}{]},
  literate={á}{{\'a}}1 {é}{{\'e}}1 {í}{{\'i}}1 {ó}{{\'o}}1 {ú}{{\'u}}1
           {ñ}{{\~n}}1 {Á}{{\'A}}1 {É}{{\'E}}1 {Í}{{\'I}}1 {Ó}{{\'O}}1
           {Ú}{{\'U}}1 {Ñ}{{\~N}}1 {¿}{{?`}}1 {¡}{{!`}}1
}
\lstdefinestyle{shell}{basicstyle=\ttfamily\fontsize{7}{8.3}\selectfont\color{codigotexto},backgroundcolor=\color{codigofondo},frame=none,xleftmargin=6pt,xrightmargin=6pt,framexleftmargin=6pt,framexrightmargin=6pt,framextopmargin=5pt,framexbottommargin=5pt,columns=flexible,showstringspaces=false}

% ---------------------- CAPTURAS OPCIONALES ---------------------------
% \captura[ancho]{archivo}{descripción}: si no existe el archivo, deja un marcador.
\newcommand{\captura}[3][0.98]{%
  \IfFileExists{#2}{%
    \includegraphics[width=#1\linewidth,height=0.48\textheight,keepaspectratio]{#2}}{%
    \colorbox{lila}{\parbox[c][0.36\textheight][c]{#1\linewidth}{\centering
      \scriptsize\color{gris}[Insertar captura]\\ #3\\ \texttt{#2}}}}}

\begin{document}

% ===================== 1. PORTADA (fondo negro) ======================
{\setbeamercolor{background canvas}{bg=tinta}
\begin{frame}[plain]
  \vspace{0.1cm}
  {\color{acentoclaro}\scriptsize\bfseries\MakeUppercase{UTEQ \ · \ Ingeniería de Requerimientos (ISR-401)}}\par
  \vspace{0.35cm}
  {\color{white}\fontsize{19}{23}\selectfont\bfseries Técnicas de elicitación y\\ herramientas automatizadas de apoyo\\ a la Ingeniería de Requerimientos\par}
  \vspace{0.25cm}
  {\color{acentoclaro}\rule{2.2cm}{2pt}}\par
  \vspace{0.2cm}
  {\color{codigotexto}\normalsize GA-SUM-02 · EXP-1 · Presentación del 1.\textsuperscript{er} corte\\
   Demostración con StrictDoc sobre el PFC Pauly Creative's Craft\par}
  \vfill
  \renewcommand{\arraystretch}{1.25}
  \begin{tabular}{@{}>{\color{acentoclaro}\tiny\bfseries}p{1.8cm}>{\color{white}\footnotesize}p{5.0cm}>{\color{acentoclaro}\tiny\bfseries}p{1.3cm}>{\color{white}\footnotesize}p{4.6cm}@{}}
    ESTUDIANTE & \Estudiante & DOCENTE & \Docente\\
    CORREO & \Correo & EQUIPO & \Equipo\\
    CARRERA & \Carrera\ · \Semestre & FECHA & \Fecha\\
  \end{tabular}
  \vspace{0.1cm}
\end{frame}}

% ===================== 2. CONTENIDO =================================
\begin{frame}{Contenido}{Agenda}
  \begin{columns}[T]
    \column{0.46\textwidth}
    {\color{acento}\large\bfseries 01}\hspace{10pt}Introducción y propósito\par\vspace{12pt}
    {\color{acento}\large\bfseries 02}\hspace{10pt}Contexto del PFC\par\vspace{12pt}
    {\color{acento}\large\bfseries 03}\hspace{10pt}Parte A · Técnicas de elicitación\par\vspace{12pt}
    {\color{acento}\large\bfseries 04}\hspace{10pt}Parte B · Herramientas de apoyo a la IR\par\vspace{12pt}
    \column{0.46\textwidth}
    {\color{acento}\large\bfseries 05}\hspace{10pt}Parte C · Recomendación para el PFC\par\vspace{12pt}
    {\color{acento}\large\bfseries 06}\hspace{10pt}Parte D · Demostración con StrictDoc\par\vspace{12pt}
    {\color{acento}\large\bfseries 07}\hspace{10pt}Conclusiones\par\vspace{12pt}
    {\color{acento}\large\bfseries 08}\hspace{10pt}Uso de IA y referencias\par\vspace{12pt}
  \end{columns}
\end{frame}

% ===================== 3. INTRODUCCIÓN ===============================
\begin{frame}{Introducción y propósito}{Introducción}
  \begin{columns}[T]
    \column{0.36\textwidth}
    \vspace{6pt}
    {\fontsize{22}{26}\selectfont\bfseries Elicitar bien,\\ y \textcolor{acento}{trazar}\\ lo elicitado.\par}
    \column{0.58\textwidth}
    \small
    \begin{enumerate}
      \item La elicitación es una actividad central de la IR: descubrir y entender las necesidades de los interesados \cite{pohl2025}.
      \item IREB propone elegir las técnicas según el interesado, la fuente de información y el conocimiento que se busca \cite{ireb2026}.
      \item Sin herramientas de apoyo, los requisitos se dispersan y se pierde la trazabilidad \cite{wiegers2013}.
    \end{enumerate}
  \end{columns}
  \nota{Propósito}{\small Describir técnicas de elicitación, comparar herramientas de apoyo a la IR y demostrar \textbf{StrictDoc} sobre mi PFC.}
\end{frame}

% ===================== 4. CONTEXTO PFC ===============================
\begin{frame}{Pauly Creative's Craft}{Contexto del PFC}
  \begin{columns}[T]
    \column{0.30\textwidth}
    \etiqueta{Organización}\small Local de manualidades en Quevedo, Los Ríos.\par\vspace{8pt}
    \etiqueta{Usuaria única}\small PI-01 (propietaria): administra el negocio, recibe pedidos, registra abonos y entregas.
    \column{0.30\textwidth}
    \etiqueta{Solución}\small Aplicación web para clientes, proveedores, productos, ventas, pagos, abonos y entregas.\par\vspace{8pt}
    \etiqueta{Foco crítico}\small Seguimiento de pedidos personalizados, de la solicitud a la entrega.
    \column{0.30\textwidth}
    \etiqueta{Técnicas aplicadas en el PFC}
    \small
    \begin{enumerate}
      \item Entrevista semiestructurada a PI-01
      \item Observación no participante a PI-01
      \item Cuestionarios a clientes
    \end{enumerate}
  \end{columns}
\end{frame}

% ===================== 5. ENTREVISTA =================================
\begin{frame}{Entrevista semiestructurada}{Parte A · Elicitación}
  \etiqueta{Pasos}
  \noindent\paso{01}{Definir objetivos y preparar un guion base}\hfill
  \paso{02}{Agendar y explicar el propósito a PI-01}\hfill
  \paso{03}{Conducir con preguntas abiertas y repreguntas}\hfill
  \paso{04}{Registrar con notas o audio, con consentimiento}\hfill
  \paso{05}{Resumir y validar lo entendido con PI-01}
  \vspace{10pt}
  \begin{columns}[T]
    \column{0.29\textwidth}
    \etiqueta{Ventajas}
    \footnotesize\begin{itemize}
      \item Profundidad y flexibilidad
      \item Aclara dudas en el momento
      \item Construye relación con la usuaria
    \end{itemize}
    \column{0.29\textwidth}
    \etiqueta{Limitaciones}
    \footnotesize\begin{itemize}
      \item Depende de la memoria y subjetividad de PI-01
      \item No capta el conocimiento tácito
      \item Exige tiempo de preparación y análisis
    \end{itemize}
    \column{0.29\textwidth}
    \etiqueta{Cuándo usarla}
    \footnotesize Al inicio, para entender el dominio y los objetivos. En el PFC: cómo PI-01 recibe pedidos y registra abonos.
  \end{columns}
\end{frame}

% ===================== 6. OBSERVACIÓN ================================
\begin{frame}{Observación no participante}{Parte A · Elicitación}
  \etiqueta{Pasos}
  \noindent\paso{01}{Definir el foco, p. ej., el ciclo de un pedido}\hfill
  \paso{02}{Obtener la autorización de PI-01}\hfill
  \paso{03}{Observar sin intervenir en su trabajo}\hfill
  \paso{04}{Registrar en una hoja de observación}\hfill
  \paso{05}{Analizar y contrastar con la entrevista}
  \vspace{10pt}
  \begin{columns}[T]
    \column{0.29\textwidth}
    \etiqueta{Ventajas}
    \footnotesize\begin{itemize}
      \item Revela el proceso real, no el declarado
      \item Descubre conocimiento tácito
      \item Detecta excepciones y retrabajo
    \end{itemize}
    \column{0.29\textwidth}
    \etiqueta{Limitaciones}
    \footnotesize\begin{itemize}
      \item La presencia puede alterar el comportamiento
      \item No explica los motivos por sí sola
      \item Interpretación subjetiva y costosa en tiempo
    \end{itemize}
    \column{0.29\textwidth}
    \etiqueta{Cuándo usarla}
    \footnotesize En procesos operativos y tácitos. En el PFC: cómo PI-01 anota hoy abonos y estados de pedidos.
  \end{columns}
\end{frame}

% ===================== 7. COMPLEMENTO ================================
\begin{frame}{Cómo se complementan}{Parte A · Elicitación}
  \small
  \renewcommand{\arraystretch}{1.2}
  \setlength{\tabcolsep}{5pt}
  \begin{tabular}{@{}p{3.0cm}p{3.7cm}p{3.1cm}p{2.7cm}@{}}
    \toprule
    {\color{acento}\bfseries Técnica} & {\color{acento}\bfseries Obtiene} & {\color{acento}\bfseries Fuente en el PFC} & {\color{acento}\bfseries Riesgo}\\
    \midrule
    Entrevista semiestructurada & Necesidades, reglas y expectativas declaradas & PI-01 & Sesgo y olvido\\
    Observación no participante & Flujo real y conocimiento tácito & PI-01 en su trabajo diario & Efecto del observador\\
    Cuestionario & Opiniones de muchos interesados & Clientes & Respuestas superficiales\\
    \bottomrule
  \end{tabular}
  \nota{Triangulación}{\small Contrastar lo dicho (entrevista), lo hecho (observación) y lo esperado por los clientes (cuestionario) reduce omisiones y contradicciones \cite{pohl2025,ireb2026}.}
\end{frame}

% ===================== 8. COMPARACIÓN ================================
\begin{frame}{Jira y Enterprise Architect frente a StrictDoc}{Parte B · Herramientas de apoyo a la IR}
  \scriptsize
  \renewcommand{\arraystretch}{1.12}
  \setlength{\tabcolsep}{4pt}
  \arrayrulecolor{gris!40}
  \begin{tabular}{@{}p{2.1cm}p{5.1cm}p{5.4cm}@{}}
    \toprule
    & {\color{acento}\bfseries Jira / Enterprise Architect} & {\color{acento}\bfseries StrictDoc}\\
    \midrule
    {\bfseries Enfoque} & Jira: incidencias y ALM. EA: modelado y repositorio de requisitos & Requisitos y especificaciones como documentos de texto\\
    {\bfseries Formato} & Repositorio propio (base de datos o modelo) & Texto plano estructurado (\texttt{.sdoc})\\
    {\bfseries Funciones} & Flujos de trabajo, tableros, UML, matrices y reportes & Edición en texto o interfaz web, exportación (HTML y otros), relaciones\\
    {\bfseries Trazabilidad} & Enlaces entre tickets o elementos; vistas según la herramienta & Relaciones padre-hijo en texto y vistas generadas automáticamente\\
    {\bfseries Versionado} & Historial propio de la herramienta & Git: ramas, \emph{diff}, historial y revisión\\
    {\bfseries Madurez} & Alta: amplia adopción y ecosistema & Menor: proyecto abierto, joven y en evolución\\
    {\bfseries Costo y curva} & Licencias o suscripción; curva media-alta & Código abierto; curva baja si se conoce Git y la terminal\\
    \bottomrule
  \end{tabular}
  \vspace{2pt}\par
  {\color{gris}\tiny Elaboración propia a partir de \cite{wiegers2013,strictdocdocs}. Verificar versiones y licencias vigentes.}
\end{frame}

% ===================== 9. STRICTDOC ==================================
\begin{frame}{¿Cómo funciona StrictDoc?}{Parte B · Herramientas de apoyo a la IR}
  \vspace{4pt}
  \begin{center}
  \begin{tikzpicture}[node distance=0.9cm,
      caja/.style={draw=tinta,thin,align=center,minimum height=1.0cm,minimum width=2.6cm,font=\scriptsize},
      fin/.style={caja,fill=acento,draw=acento,text=white},
      flecha/.style={-{Stealth[length=2mm]},thin,color=tinta}]
    \node[caja] (a) {Archivos \texttt{.sdoc}\\texto plano};
    \node[caja,right=of a] (b) {Git / GitHub\\historial y revisión};
    \node[caja,right=of b] (c) {\texttt{strictdoc}\\export / server};
    \node[fin,right=of c] (d) {HTML y vistas\\de trazabilidad};
    \draw[flecha] (a)--(b); \draw[flecha] (b)--(c); \draw[flecha] (c)--(d);
  \end{tikzpicture}
  \end{center}
  \vspace{6pt}
  \begin{columns}[T]
    \column{0.30\textwidth}
    {\color{acento}\large\bfseries 01}\par\small Cada elemento tiene \texttt{UID}, título, enunciado y relaciones \cite{strictdocdocs}.
    \column{0.30\textwidth}
    {\color{acento}\large\bfseries 02}\par\small Los cambios se revisan como código: \emph{commit}, \emph{diff} y \emph{pull request} \cite{strictdocrepo}.
    \column{0.30\textwidth}
    {\color{acento}\large\bfseries 03}\par\small La trazabilidad se genera al exportar, sin matrices manuales.
  \end{columns}
\end{frame}

% ===================== 10. RECOMENDACIÓN =============================
\begin{frame}{Recomendación para el PFC}{Parte C · Recomendación}
  {\fontsize{16}{20}\selectfont\bfseries Usar \textcolor{acento}{StrictDoc} como herramienta de apoyo a la IR de Pauly Creative's Craft.\par}
  \vspace{6pt}
  \begin{columns}[T]
    \column{0.47\textwidth}
    \etiqueta{Por qué}
    \small\begin{itemize}
      \item Un equipo y una usuaria: Jira o EA serían sobredimensionados
      \item Costo cero; el PFC ya usa Git y GitHub
      \item Trazabilidad N → RF → CP verificable y versionada
      \item Requisitos exportables y revisables por el docente
    \end{itemize}
    \column{0.47\textwidth}
    \etiqueta{Riesgos y mitigación}
    \small\begin{itemize}
      \item PI-01 no es técnica: no edita \texttt{.sdoc}; valida la vista HTML exportada
      \item Menor madurez: fijar la versión instalada y respaldar en GitHub
      \item Curva en terminal: guía breve de comandos en el repositorio
    \end{itemize}
  \end{columns}
\end{frame}

% ===================== 11. DEMO 1 ====================================
\begin{frame}[fragile]{N-01 y RF-01 en formato \texttt{.sdoc}}{Parte D · Demostración}
  \begin{columns}[T]
    \column{0.45\textwidth}
    \etiqueta{Necesidad}
\begin{lstlisting}[style=sdoc]
[DOCUMENT]
TITLE: Pauly Creative's Craft

[REQUIREMENT]
UID: N-01
TAGS: Necesidad
TITLE: Seguimiento de pedidos
STATEMENT: >>>
PI-01 necesita conocer el estado
de cada pedido personalizado,
desde la solicitud hasta la
entrega.
<<<
\end{lstlisting}
    \column{0.47\textwidth}
    \etiqueta{Requisito funcional}
\begin{lstlisting}[style=sdoc]
[REQUIREMENT]
UID: RF-01
TAGS: Funcional
TITLE: Registrar pedido
STATEMENT: >>>
El sistema debe permitir a PI-01
registrar un pedido personalizado
con cliente, descripción, fecha de
entrega, precio y estado inicial.
<<<
RELATIONS:
- TYPE: Parent
  VALUE: N-01
\end{lstlisting}
  \end{columns}
\end{frame}

% ===================== 12. DEMO 2 ====================================
\begin{frame}[fragile]{CP-01 y comandos}{Parte D · Demostración}
  \begin{columns}[T]
    \column{0.58\textwidth}
\begin{lstlisting}[style=sdoc]
[REQUIREMENT]
UID: CP-01
TAGS: Prueba
TITLE: Registrar pedido con datos válidos
STATEMENT: >>>
Dado que PI-01 ingresa cliente, descripción, fecha y precio
válidos, cuando guarda el pedido, entonces el sistema lo
registra con estado "Solicitado".
<<<
RELATIONS:
- TYPE: Parent
  VALUE: RF-01
\end{lstlisting}
    \column{0.36\textwidth}
    \etiqueta{Comandos}
\begin{lstlisting}[style=shell]
pip install strictdoc
strictdoc --version
strictdoc export .
strictdoc server .
\end{lstlisting}
    \vspace{8pt}
    \etiqueta{Cadena de trazabilidad}
    {\small N-01 $\leftarrow$ RF-01 $\leftarrow$ CP-01\par}
  \end{columns}
\end{frame}

% ===================== 13. CAPTURAS ==================================
\begin{frame}{Resultados en StrictDoc}{Parte D · Demostración}
	\begin{columns}[T]
		\column{0.47\textwidth}
		\centering
		\captura{evidencias/captura-export-html.png}{documento exportado con N-01, RF-01 y CP-01}\par
		{\color{gris}\scriptsize Figura 1. Documento exportado en HTML}
		\column{0.47\textwidth}
		\centering
		\captura{evidencias/captura-trazabilidad.png}{vista de trazabilidad N-01, RF-01, CP-01}\par
		{\color{gris}\scriptsize Figura 2. Vista de trazabilidad}
	\end{columns}
	\vspace{8pt}
	{\color{gris}\scriptsize Repositorio: \texttt{\RepoURL}}
\end{frame}

% ===================== 14. CONCLUSIONES ==============================
\begin{frame}{Conclusiones}{Cierre}
  \vspace{4pt}
  \begin{columns}[T]
    \column{0.45\textwidth}
    \conc{01}{Entrevista y observación se complementan: lo declarado frente a lo realmente hecho.}
    \conc{02}{Los cuestionarios a clientes amplían el alcance de los requisitos.}
    \column{0.45\textwidth}
    \conc{03}{Las herramientas CASE aportan madurez; StrictDoc aporta ligereza, versionado y trazabilidad automática.}
    \conc{04}{Para Pauly Creative's Craft, StrictDoc es adecuada, mitigando sus riesgos.}
  \end{columns}
\end{frame}

% ===================== 15. USO DE IA =================================
\begin{frame}{Declaración de uso de IA como apoyo}{Normativa UTEQ}
  \begin{columns}[T]
    \column{0.44\textwidth}
    \etiqueta{Herramienta}\small Claude (Anthropic), interfaz web.\par\vspace{8pt}
    \etiqueta{Versión}\small Claude Sonnet 5.5, noviembre de 2026.\par\vspace{8pt}
    \etiqueta{Propósito}\small Apoyo en la estructura de la presentación, la redacción inicial y el código LaTeX/Beamer.
    \column{0.46\textwidth}
    \etiqueta{Verificación}\small Contraste de las fuentes bibliográficas, prueba real de StrictDoc y revisión del contenido frente al PFC por el estudiante.\par\vspace{8pt}
    \etiqueta{Responsabilidad}\small El contenido final es responsabilidad del estudiante.
  \end{columns}
\end{frame}

% ===================== 16. REFERENCIAS ===============================
\begin{frame}[allowframebreaks]{Referencias}{Formato IEEE}
  \footnotesize
  \bibliographystyle{IEEEtran}
  \bibliography{referencias}
\end{frame}

\end{document}
